\documentclass{article}

%Basics
\usepackage{amsmath, amssymb, amsthm}
\usepackage{enumitem}
\usepackage[margin=1in]{geometry}

%Links
\usepackage[colorlinks]{hyperref}

%Colored Strikethrough
\usepackage{soul}
\setstcolor{red}

%Coloring Digits
\usepackage{xcolor}

%Formatting and Spacing
\setitemize[1]{noitemsep, parsep = 5pt, topsep = 5pt}
\setenumerate[1]{label = (\alph*), parsep = 1pt, topsep = 5pt}
\setlength\parindent{0pt}
\linespread{1.1}

%Easy Numbering
\newcounter{counter}
\newcommand{\Exercise}{Exercise \stepcounter{counter}\thecounter}

%Custom Title Fields
\newcommand{\hmwkTitle}{Exam 2 Extra Credit}
\newcommand{\hmwkDueDate}{Sunday, April 24, 2022}
\newcommand{\hmwkClass}{Honors Discrete Mathematics}
\newcommand{\hmwkClassInstructor}{Professor Gerandy Brito}
\newcommand{\hmwkSection}{Spring 2022.}
\newcommand{\hmwkAuthorName}{Sarthak Mohanty}

%Headers and Footers
\usepackage{fancyhdr}
\usepackage{extramarks}
\pagestyle{fancy}
\lhead{CS 2051}
\chead{\hmwkClass \ (\hmwkClassInstructor)}
\rhead{Exam 2 EC}
\cfoot{\thepage}
\renewcommand\headrulewidth{0.4pt}
\renewcommand\footrulewidth{0.4pt}

\title{
    \vspace{2in}
    \textbf{\hmwkClass:\\ \hmwkTitle} \\
    \normalsize\vspace{0.1in}\small{Due\ on\ \hmwkDueDate\ at 11:59pm} \\
    \vspace{0.1in}\large{\textit{\hmwkClassInstructor\ \hmwkSection}} \\
    \vspace{1in}
    \begin{minipage}[t]{0.5\columnwidth}
    {\footnotesize \textbf{IMPORTANT:} For this assignment, you may NOT collaborate with other students, nor use resources outside the ones from class. Think of it as a take-home exam that is open book, open notes.}
    \end{minipage}
    \vspace{1in}
    \author{\textbf{\hmwkAuthorName}}
    \date{}
}

\begin{document}
    
\maketitle
\pagebreak

\subsection*{\Exercise \ {\bf [1 pt]}}
    Every time you visited The Farm, you went to the barn to get some fresh milk. To milk the cows, the farmer brings in $23$ distinct cows from the pasture into the barnyard and lines them up.

    \vspace{1.5mm}
    You notice that two adjacent cows never both give milk on the same day, but no more that two adjacent cows ever give no milk on the same day. On a given day, how many different ways can the farmer get milk? (An example would be ``Cow 1 gives milk, Cow 2 does not, Cow 3 does not, etc.")

\subsection*{\Exercise \ {\bf [1 pt]}}
    Consider the grid shown below in Figure \ref{E2}.
    
    \vspace{1mm}
    \setlength{\tabcolsep}{2.7pt}
    \begin{figure}[hbp]
        \centering
        \begin{tabular}{ccccccccccccc}
            & & & & & & S \\
            & & & & & S & A & S \\
            & & & & S & A & R & A & S \\
            & & & S & A & R & T & R & A & S \\
            & & S & A & R & T & H & T & R & A & S \\
            & S & A & R & T & H & A & H & T & R & A & S \\
            S & A & R & T & H & A & K & A & H & T & R & A & S \\
            & S & A & R & T & H & A & H & T & R & A & S \\
            & & S & A & R & T & H & T & R & A & S \\
            & & & S & A & R & T & R & A & S \\
            & & & & S & A & R & A & S \\
            & & & & & S & A & S \\
            & & & & & & S
        \end{tabular}
        \caption{A diamond-shaped, narcissistic grid of letters.}
        \label{E2}
    \end{figure}
    
    Suppose you start at any $S$, and can move only left, right, down, or up to adjacent letters. In how many ways can the palindrome SARTHAKAHTRAS be read if
    \begin{enumerate}
        \item the same letter can be used more than once in each sequence?
        \item \st{the same letter cannot be used more than one in each sequence?}
    \end{enumerate}

\subsection*{\Exercise \ {\bf [1 pt, Challenge]}}
    Let $\mathcal{P}$ be the set of all prime numbers. Let $M$ be a subset of $\mathcal{P}$ with at least $3$ elements. Choose any proper subset $A$ of $M$. Consider the number $$n_{A} = -1 + \prod_{p \in A} p.$$ Suppose that any prime divisor of $n_{A}$ lies in $M$ for all $A \subseteq M$. Show that $M = \mathcal{P}$.

\subsection*{\Exercise \ [0 pts, Fall '21 Challenge]}
    Let $N$ be a natural number, written in base 10. Show that there is a power of two such that its first digits are exactly those in $N$, in the same order. For example,
    
    \vspace{1.5mm}
    for $N = \textcolor{red}{3}$ we have $2^{5} = \textcolor{red}{3}2$; \\
    for $N = \textcolor{blue}{12}$ we have $2^{7} = \textcolor{blue}{12}8$; \\
    for $N = \textcolor{cyan}{102}$ we have $2^{10} = \textcolor{cyan}{102}4$.

\end{document}